\section{Theoretical Findings: Research Workflow and Proofs}
\label{app:proof-status}

This appendix first describes the research design and workflow used to
develop and check the findings in Section~\ref{sec:our-results}.
Complete written proofs follow in the order of the three main-text cases.

\subsection{Research design and workflow}
\label{app:research-method}
We supplied advanced reasoning models, including GPT-6 Astra with Ultra
reasoning and GPT-6 Pro, with the Lean library, paper maps, and relevant source
manuscripts. The models compared definitions and assumptions across papers,
looked for connections and simpler arguments, and proposed new questions.
Promising directions were developed through successive rounds of statement
refinement, proof construction, example testing, and review. The number of
iterations and the supplied material varied with the question.

The workflow combined mathematical exploration with focused checks. Follow-up
sessions examined proofs, searched for counterexamples, ran finite checks, and
formalized selected statements in Lean. For each featured finding,
the following proof subsections identify its source material, additional
mathematical contribution, written proof, and Lean verification status.



\paragraph{Overview of proofs and verification.}
Table~\ref{tab:proof-guide} locates each proof and distinguishes the
written arguments and their Lean verification.

\begin{table}[ht]
\centering
\caption{Proofs of the featured findings and scope of formal verification.}
\label{tab:proof-guide}
\small
\begin{tabular}{@{}>{\raggedright\arraybackslash}p{0.30\linewidth}>{\raggedright\arraybackslash}p{0.15\linewidth}>{\raggedright\arraybackslash}p{0.47\linewidth}@{}}
\toprule
Finding & Proof & Lean scope \\
\midrule
Section~\ref{sec:precision}: local invariant repair & Appendix~\ref{app:pareto-repair} & Counterexample, corrected score bound, and the repaired generation guarantee, including repeated observations. \\
Theorem~\ref{thm:proper-replay}: finite proper replay & Appendix~\ref{app:replay} & Three-language impossibility, both directions of the characterization, and finite-state consequences. \\
Theorem~\ref{thm:staircase-frontier}: exact staircase frontier & Appendix~\ref{app:p29-frontier} & Exact guarantees, lower bounds for arbitrary deterministic generators, and Pareto optimality. \\
\bottomrule
\end{tabular}
\end{table}

\paragraph{Proof materials.}
The complete Lean sources, theorem-to-code correspondence, and reproducible
verification commands accompany these written proofs in the
anonymous artifact repository.
The formalization concerns the deterministic mathematical models stated
below; it does not assert effective oracle implementations or running-time
bounds.

\subsection{Identifying and repairing a proof gap in Pareto-optimal generation}
\label{app:pareto-repair}

In this subsection we give a repair to the proof of the generation-time
theorem of \citet[][P13]{charikar2026pareto}. During formalization we
found that published Claim~7 (arXiv Claim~3.2), used in published
Theorem~8 (arXiv Theorem~4), can fail when Procedure~1 stores an empty
subcollection to indicate that no candidate exists. We explain the
claim's role, exhibit the failure, and prove the numerical bound needed
to complete the generation argument. The repaired proof preserves the
original algorithm and its generation-time guarantee. References below
use the published ALT~2026 numbering.

\subsubsection{The insertion procedure and the original claim}

Let $L_1,L_2,\ldots$ be infinite languages over a countable universe $U$.
Procedure~1 inserts these languages one at a time into a finite ordered
list. It assigns each inserted language $L_i$ a nonnegative integer
$m_i$ and a stored subcollection $C_i$; both remain fixed after that
insertion. Here $m_i$ denotes the source paper's $m^*(L_i)$. The order,
and hence the languages preceding $L_i$, can change when later languages
are inserted.

Identify a subcollection with its set of indices. For a finite
index set $S$, write
\[
 I(D)=\bigcap_{j\in D}L_j,
 \qquad
 \mathcal C_i(S)=
 \{D\subseteq S:i\in D\text{ and }I(D)\text{ is finite}\}.
\]
Thus a member of $\mathcal C_i(S)$ is a subcollection containing $L_i$
whose intersection is finite. Define
\[
 M_i(S)=\max\bigl(\{0\}\cup
       \{|I(D)|:D\in\mathcal C_i(S)\}\bigr).
\]
The maximum exists because $S$ has finitely many subcollections. If
$\mathcal C_i(S)$ is nonempty, it is the largest finite intersection
size attained by a candidate. If the candidate set is empty, its value
is zero by convention. In particular, $M_i(S)=0$ does not distinguish
an empty candidate set from a nonempty candidate set whose intersections
are all empty.

Suppose a new language $L_q$ is immediately to the right of $L_i$ during
its insertion, and let $S_i$ consist of $i$ and all indices preceding
$i$ in the old order. The procedure computes
\[
 m_{\mathrm{check}}=M_q(S_i\cup\{q\})
\]
and moves $q$ to the left of $i$ if
$m_{\mathrm{check}}\le m_i$. It repeats this comparison while moving
leftward. The checked maximum is recomputed in the relevant prefix;
it need not equal the score eventually assigned to $q$. Once the
insertion ends, let $S_q$ be the prefix through $q$. The procedure sets
$m_q=M_q(S_q)$ and chooses $C_q\in\mathcal C_q(S_q)$ attaining this
maximum if a candidate exists. Otherwise it sets $C_q=\varnothing$
and $m_q=0$.

Claim~7 asserts that the stored subcollection for an old language
continues to be a maximizing candidate in its current prefix. In the
notation above, its maximizing-witness assertion is
\[
 C_i\in\mathop{\rm arg\,max}_{D\in\mathcal C_i(S_i)} |I(D)|.
\]
The proof of Theorem~8 uses this assertion to bound the size of a finite
intersection involving $L_i$. Specifically, the generator scans the
ordered languages, retaining a language when it contains every observed
point and its addition leaves the retained languages with an infinite
intersection. To show that the true target $L_i$ is retained after more
than $m_i$ distinct observations, the proof needs the implication
\[
 D\in\mathcal C_i(S_i)\quad\Longrightarrow\quad |I(D)|\le m_i.
\]
The assertion that the particular stored subcollection $C_i$ attains
the maximum is stronger than this implication.

\subsubsection{Failure of the maximizing-witness assertion}

Let $A$ and $B$ be disjoint infinite languages, inserted in that order.
At the first insertion, the only subcollection containing $A$ is
$\{A\}$, whose intersection is infinite. There is therefore no
candidate, and the procedure stores $C_A=\varnothing$ and $m_A=0$.

When $B$ is inserted, the subcollection $\{A,B\}$ has empty
intersection. It is a candidate for $B$ with score zero, whereas
$\{B\}$ still has infinite intersection. Consequently the checked
maximum is zero. The comparison $0\le m_A=0$ moves $B$ before $A$,
giving the order $(B,A)$. The prefix through $A$ now contains both
languages, and $\{A,B\}$ is a candidate for $A$ attaining the
maximum zero. However, the stored subcollection for $A$ is still
$\varnothing$. It does not contain $A$, so it is not a candidate and
cannot belong to the displayed argmax.

The failure is the identification of an empty stored subcollection
with a maximizing candidate. The empty subcollection contains no
languages; the candidate $\{A,B\}$ contains two languages whose
intersection contains no points. Assigning score zero in both cases
does not make the two subcollections interchangeable. The numerical
inequality needed by Theorem~8 remains true in this execution: every
candidate in the prefix through $A$ has intersection size at most
$m_A=0$.
The published induction controls the maximum value, but this alone
does not ensure that the stored subcollection is still a candidate.

\subsubsection{The corrected statement and its preservation}

The repair replaces the assertion about a stored maximizing
subcollection by an upper bound on every candidate's intersection
size. The following lemma provides the induction step: it shows
that the numerical bound persists when a new language enters an
old language's prefix.

\begin{supplem}[Preservation of the score bound]
\label{lem:pareto-score-invariant}
Let $S$ be a finite set of indices and let $b\in\mathbb N$.
Suppose that every $D\in\mathcal C_i(S)$ satisfies $|I(D)|\le b$.
If $q\notin S$ and $M_q(S\cup\{q\})\le b$, then every
$D\in\mathcal C_i(S\cup\{q\})$ also satisfies $|I(D)|\le b$.
\end{supplem}
\begin{proof}
Fix $D\in\mathcal C_i(S\cup\{q\})$. There are two cases.
If $q\notin D$, then $D\subseteq S$, and it still contains $i$
and has finite intersection. Thus $D\in\mathcal C_i(S)$, so the
assumed bound gives $|I(D)|\le b$.

If $q\in D$, then $D\subseteq S\cup\{q\}$ contains $q$ and has
finite intersection. It is therefore also a candidate in the
optimization defining $M_q(S\cup\{q\})$. Hence
\[
 |I(D)|\le M_q(S\cup\{q\})\le b.
\]
Both cases establish the required inequality. The argument includes
$b=0$ and does not require a stored maximizing subcollection.
\end{proof}

\begin{suppprop}[Corrected form of Claim~7]
At every completed insertion stage, let $L_i$ be an inserted language,
let $m_i$ be its assigned score, and let $S_i$ be its current prefix,
including $i$. Then
\begin{equation}
\label{app:pareto-corrected-bound}
 \text{for every }D\subseteq S_i,\qquad
 i\in D\text{ and }I(D)\text{ finite}
 \quad\Longrightarrow\quad |I(D)|\le m_i.
\end{equation}
\end{suppprop}
\begin{proof}
We induct on the number of insertions. The assertion is vacuous
before any language is inserted. Suppose it holds for the old
order and insert $L_q$.
For the new language, its assigned score is $m_q=M_q(S_q)$, so the
bound follows directly from the definition of the maximum in its
final prefix $S_q$. In particular, it is vacuous when there is no
candidate and remains valid when the maximum is zero.

Consider an old language $L_i$. If $q$ does not move to the left of
$i$, its prefix and assigned score remain unchanged, and the
inductive hypothesis applies. If $q$ moves to the left of $i$, let
$S_i$ be its old prefix. The relative order of all old languages
is unchanged, so its new prefix is exactly $S_i\cup\{q\}$.
When $q$ crosses $i$, the procedure's comparison gives
\[
 M_q(S_i\cup\{q\})\le m_i.
\]
The inductive hypothesis bounds every $i$-candidate in $S_i$ by
$m_i$. Applying Lemma~\ref{lem:pareto-score-invariant} with
$b=m_i$ proves the bound in its new prefix. Subsequent movements
of $q$ farther to the left change neither this set of indices nor
$m_i$. This covers every old language and completes the induction.
\end{proof}

\subsubsection{Recovery of the generation-time guarantee}

For completeness, we give the step connecting the corrected bound
to the original generation algorithm. Let $H\subseteq L_i$ be the
finite set of distinct observations from the target $L_i$, and suppose
that $i$ is included in the finite list currently being scanned and
that $|H|\ge m_i+1$. Let $B$ be the indices retained before the
scan reaches $i$. Every retained language contains $H$, and every
index in $B$ precedes $i$. Consequently
\[
 B\cup\{i\}\subseteq S_i,
 \qquad
 H\subseteq I(B\cup\{i\}).
\]
The target satisfies the sample-consistency test because
$H\subseteq L_i$. If the scan rejected it, its intersection with
the previously retained languages would be finite. Then
$B\cup\{i\}$ would be an $i$-candidate, and
\eqref{app:pareto-corrected-bound} would imply
\[
 m_i+1\le |H|
       \le |I(B\cup\{i\})|
       \le m_i,
\]
a contradiction. Thus the scan retains $L_i$.

The scan never removes a retained language and only adds a language
when the common intersection remains infinite. Its final
intersection is therefore infinite and, because it includes $L_i$
among the intersected languages, is a subset of $L_i$. Removing
the finite observed set leaves a point in $L_i\setminus H$.
The original generator's choice of such a point is consequently a
fresh target output.

In Theorem~8, the generator considers the first $f(r)$ languages
after $r$ observations, where $f$ is nondecreasing and unbounded.
Observations may repeat; let $H_r$ be the set of distinct points
observed by round $r$. Write $g(i)=\min\{r:f(r)\ge i\}$ for
the first round at which $L_i$ is included. If
\[
 |H_r|\ge t_i:=\max\{g(i),m_i+1\},
\]
then $r\ge |H_r|\ge g(i)$, so $f(r)\ge i$, and $|H_r|>m_i$.
The target is therefore under consideration and the preceding
argument guarantees a fresh target output. This proves the same
bound in terms of distinct observations as in the source paper,
including presentations with repetitions.
The score assignments, insertion comparisons, greedy
scan, and output rule all remain unchanged. Only the intermediate
proof assertion is replaced.

\paragraph{Formal verification.}
GenLimitLib checks the two-language counterexample, the corrected
score bound through successive insertions, target retention, and the
resulting generation-time guarantee. The accompanying Lean development
also proves the version above with repeated observations, retaining the
original scheduler based on the total number of rounds. The repository
maps these statements to P13's existing modules and the additional proof
in \nolinkurl{Section4/InvariantRepair.lean}.
